% !PRINTED 467
\weberpage{467}

%<<CONT
so sind $A_1=0$, $A_2=0$, $A_3=0$ die Gleichungen von dreien dieser Axen. Durch die Substitution $\chi$ gehen $A_1, A_2, A_3$ cyklisch in einander über. Durch die Substitution $\Theta$ erleiden die Functionen $A_1, A_2, A_3$ folgende Vertauschung:

\[
\begin{pmatrix} A_1, & A_2, & A_3 \\ A_1, & -A_3, & A_2 \end{pmatrix},
\]

wie eine einfache Rechnung zeigt, wenn man die Substitution $\Theta$ wirklich ausführt und die Formeln des \S.\,115 benutzt.

Um den Gang der Rechnung wenigstens anzudeuten, sei bemerkt, dass $A_1$ durch die Substitution $\Theta$ [\S.\,115, (10)] in

\[
\alpha(\gamma^2\varepsilon^2 + \beta\gamma\varepsilon^3 + \beta^2\varepsilon^5)x_1 + \gamma(\alpha^2\varepsilon^6 + \beta^2 + \alpha\beta\varepsilon^2)x_2 + \beta(\alpha\gamma + \gamma^2\varepsilon + \alpha^2\varepsilon^3)x_3
\]

übergeht. Es ist aber nach \S.\,115, (8), (17), (11):

\[
\begin{aligned}
\gamma^2\varepsilon^2 + \beta\gamma\varepsilon^3 + \beta^2\varepsilon^5 &= \gamma^2\varepsilon^2 + \alpha^2\varepsilon^3 + \beta^2\varepsilon^5 + \alpha\varepsilon^3 \\
&= \alpha(\varepsilon^4+\varepsilon^{-4}) = h(\varepsilon-\varepsilon^{-1}) = \gamma,
\end{aligned}
\]

und daher wird der Coëfficient von $x_1$ gleich $\alpha\gamma$, wie in $A_1$, und ebenso formt man die übrigen Ausdrücke um. Da sich nun $\omega$ aus $\Theta$ und $\chi$ zusammensetzen lässt (\S.\,72), so folgt, dass die Grössen $A_1^2, A_2^2, A_3^2$ durch die Gruppe $G_{24}$ nur unter einander permutirt werden, und dass demnach ihre symmetrischen Functionen Wurzeln von Resolventen 7\textsuperscript{ten} Grades sind.

Die einfachste symmetrische Function dieser Grössen ist die Summe

\weq{4}{A_1^2 + A_2^2 + A_3^2,}

und diese wollen wir, mit einem geeigneten numerischen Factor multiplicirt, als die Unbekannte der Resolvente einführen.

Ordnen wir die Summe (4) nach $x_1, x_2, x_3$, so erhalten wir dafür einen Ausdruck von der Form

\[
\lambda(x_1^2+x_2^2+x_3^2) + \mu(x_2x_3+x_3x_1+x_1x_2),
\]

worin nach (3):

\[
\begin{aligned}
\lambda &= (\alpha^2\gamma^2\varepsilon^{-1} + \beta^2\gamma^2\varepsilon^{-2} + \alpha^2\beta^2\varepsilon^3)\varepsilon \\
\mu &= 2\alpha\beta\gamma(\alpha\varepsilon+\beta\varepsilon^{-3}+\gamma\varepsilon^2)\varepsilon.
\end{aligned}
\]

Hierbei ist der Factor $\varepsilon$ aus der Klammer gezogen, damit der andere Factor durch die Vertauschung von $\varepsilon$ und $\varepsilon^2$ ungeändert bleibt, und sich daher rational durch $\sqrt{-7}$ ausdrücken l\"asst (Bd. I, \S.\,171).
